% Source PDF: source/普通高考/2015/2015天津理.pdf, page 1 (question 3).
% PDF embedded-bitmap attempt: pdfimages -list returned no image objects; the flowchart is vector line art.
% Grid evidence: tmp/review_2015_tj_science/grid/page_grid100.png and q3_grid50.png,
% positioned from the ungridded page render; comparison uses q3_orig_fig.png.
% Node -> directed-edge list: 开始 -> (S=20,i=1) -> (i=2i) -> (S=S-i) -> (i>5?);
% i>5? 的 是 -> 输出 S -> 结束；否 -> right loop-back -> (i=2i).
% solid segments: all node outlines and connectors; dashed segments: none.
% labels: branch labels are placed beside the corresponding diamond exits; every node is
% locally zero-indented and centered with compact padding.

\begin{tikzpicture}[
  x=1cm,
  y=1cm,
  >=Stealth,
  line width=0.45pt,
  line cap=round,
  line join=round,
  font=\normalsize,
  flowtext/.style={anchor=center, align=center,
    execute at begin node={\setlength{\parindent}{0pt}}},
  flowbox/.style={draw, minimum height=0.68cm, inner xsep=4pt, inner ysep=2pt,
    align=center, execute at begin node={\setlength{\parindent}{0pt}}},
  terminal/.style={flowbox, rounded corners=0.18cm, minimum width=1.15cm},
  io/.style={flowbox, trapezium, trapezium left angle=72, trapezium right angle=108,
    minimum height=0.72cm},
  decision/.style={flowbox, diamond, aspect=2.9, minimum width=3.25cm,
    minimum height=0.92cm, inner xsep=1pt, inner ysep=1pt},
]
  % Main top-to-bottom sequence.
  \node[terminal] (start) at (0,8.75) {开始};
  \node[flowbox, minimum width=3.40cm] (init) at (0,7.45) {$S=20,i=1$};
  \node[flowbox, minimum width=1.55cm] (double) at (0,6.05) {$i=2i$};
  \node[flowbox, minimum width=2.50cm] (subtract) at (0,4.65) {$S=S-i$};
  \node[decision] (test) at (0,3.25) {$i>5?$};
  \node[io, minimum width=1.95cm] (output) at (0,1.85) {输出 $S$};
  \node[terminal] (finish) at (0,0.55) {结束};

  \draw[->] (start.south) -- (init.north);
  % The initial trunk and the loop return share one explicit entry into i=2i;
  % place the return arrowhead before the merge and leave an unmarked tail.
  \coordinate (loop-merge) at (0,6.75);
  \coordinate (loop-tip) at (0.38,6.75);
  \draw (init.south) -- (loop-merge);
  \draw[->] (loop-merge) -- (double.north);
  \draw[->] (double.south) -- (subtract.north);
  \draw[->] (subtract.south) -- (test.north);
  \draw[->] (test.south) -- node[flowtext, right=2pt] {是} (output.north);
  \draw[->] (output.south) -- (finish.north);

  % 否 branch: the independent right-side loop returns to the i=2i entry.
  \draw (test.east) -- (2.85,3.25) -- (2.85,6.75);
  \draw[->] (2.85,6.75) -- (loop-tip);
  \draw (loop-tip) -- (loop-merge);
  \node[flowtext, anchor=west, inner sep=1pt] at (2.08,3.02) {否};
\end{tikzpicture}
